% Options for packages loaded elsewhere
\PassOptionsToPackage{unicode}{hyperref}
\PassOptionsToPackage{hyphens}{url}
\PassOptionsToPackage{dvipsnames,svgnames,x11names}{xcolor}
\documentclass[
  11pt,
  a4paper,
]{article}
\usepackage{xcolor}
\usepackage[margin=1in]{geometry}
\usepackage{amsmath,amssymb}
\setcounter{secnumdepth}{-\maxdimen} % remove section numbering
\usepackage{iftex}
\ifPDFTeX
  \usepackage[T1]{fontenc}
  \usepackage[utf8]{inputenc}
  \usepackage{textcomp} % provide euro and other symbols
\else % if luatex or xetex
  \usepackage{unicode-math} % this also loads fontspec
  \defaultfontfeatures{Scale=MatchLowercase}
  \defaultfontfeatures[\rmfamily]{Ligatures=TeX,Scale=1}
\fi
\usepackage{lmodern}
\ifPDFTeX\else
  % xetex/luatex font selection
  \ifLuaTeX
    \usepackage{luaotfload}
    \directlua{luaotfload.add_fallback("mainfontfallback",{
      "NotoSerifBengali:mode=harf"
    })}
  \fi
  \setmainfont[,RawFeature={fallback=mainfontfallback}]{Latin Modern Roman}
  \ifLuaTeX
    \usepackage{luaotfload}
    \directlua{luaotfload.add_fallback("sansfontfallback",{
      "NotoSansBengali:mode=harf"
    })}
  \fi
  \setsansfont[,RawFeature={fallback=sansfontfallback}]{Latin Modern Sans}
  \ifLuaTeX
    \usepackage{luaotfload}
    \directlua{luaotfload.add_fallback("monofontfallback",{
      "NotoSansBengali:mode=harf"
    })}
  \fi
  \setmonofont[,RawFeature={fallback=monofontfallback}]{DejaVu Sans Mono}
\fi
% Use upquote if available, for straight quotes in verbatim environments
\IfFileExists{upquote.sty}{\usepackage{upquote}}{}
\IfFileExists{microtype.sty}{% use microtype if available
  \usepackage[]{microtype}
  \UseMicrotypeSet[protrusion]{basicmath} % disable protrusion for tt fonts
}{}
\makeatletter
\@ifundefined{KOMAClassName}{% if non-KOMA class
  \IfFileExists{parskip.sty}{%
    \usepackage{parskip}
  }{% else
    \setlength{\parindent}{0pt}
    \setlength{\parskip}{6pt plus 2pt minus 1pt}}
}{% if KOMA class
  \KOMAoptions{parskip=half}}
\makeatother
\setlength{\emergencystretch}{3em} % prevent overfull lines
\providecommand{\tightlist}{%
  \setlength{\itemsep}{0pt}\setlength{\parskip}{0pt}}
% paper_preamble.tex — shared LaTeX preamble for the PCMA/SI paper set.
%
% Compiled with lualatex (NOT pdflatex) because several papers contain
% Bengali script: the CLL paper realises one cognitive state into
% English, Bengali and Banglish, so the Bengali output is EVIDENCE and
% cannot be transliterated away without destroying the point being made.
%
% Font selection is NOT done here. pandoc >= 3 has native
% `mainfontfallback` / `monofontfallback` variables that emit the
% luaotfload fallback registration in the correct position relative to
% unicode-math; doing it by hand in this header lands after unicode-math
% and breaks font loading outright. See tools/build_papers.sh.

\usepackage{amsmath}
\usepackage{amssymb}
\usepackage{mathtools}   % \xleftrightarrow etc. used in cll.md §5
\usepackage{booktabs}

% ---------------------------------------------------------------------
% Unicode symbols appearing in PROSE rather than inside $...$.
%
% The papers were written as markdown and read as markdown, so they use
% real characters — "the ω-boost", "α-child", "ΔR coherence" — instead of
% LaTeX macros. Under unicode-math those codepoints are math-mode-only
% and abort the build ("\symbballowed only in math mode"). Mapping each
% to \ensuremath keeps the markdown readable on GitHub AND compilable,
% which is the whole point of keeping one source.
% ---------------------------------------------------------------------
\usepackage{newunicodechar}

% Greek, lower case
\newunicodechar{α}{\ensuremath{\alpha}}
\newunicodechar{β}{\ensuremath{\beta}}
\newunicodechar{γ}{\ensuremath{\gamma}}
\newunicodechar{δ}{\ensuremath{\delta}}
\newunicodechar{ε}{\ensuremath{\varepsilon}}
\newunicodechar{η}{\ensuremath{\eta}}
\newunicodechar{θ}{\ensuremath{\theta}}
\newunicodechar{κ}{\ensuremath{\kappa}}
\newunicodechar{λ}{\ensuremath{\lambda}}
\newunicodechar{μ}{\ensuremath{\mu}}
\newunicodechar{π}{\ensuremath{\pi}}
\newunicodechar{ρ}{\ensuremath{\rho}}
\newunicodechar{σ}{\ensuremath{\sigma}}
\newunicodechar{τ}{\ensuremath{\tau}}
\newunicodechar{φ}{\ensuremath{\varphi}}
\newunicodechar{χ}{\ensuremath{\chi}}
\newunicodechar{ψ}{\ensuremath{\psi}}
\newunicodechar{ω}{\ensuremath{\omega}}
\newunicodechar{ν}{\ensuremath{\nu}}
\newunicodechar{ξ}{\ensuremath{\xi}}
\newunicodechar{ζ}{\ensuremath{\zeta}}
\newunicodechar{ι}{\ensuremath{\iota}}

% Greek, upper case
\newunicodechar{Δ}{\ensuremath{\Delta}}
\newunicodechar{Φ}{\ensuremath{\Phi}}
\newunicodechar{Σ}{\ensuremath{\Sigma}}
\newunicodechar{Ψ}{\ensuremath{\Psi}}
\newunicodechar{Ω}{\ensuremath{\Omega}}
\newunicodechar{Γ}{\ensuremath{\Gamma}}
\newunicodechar{Λ}{\ensuremath{\Lambda}}
\newunicodechar{Θ}{\ensuremath{\Theta}}
\newunicodechar{Π}{\ensuremath{\Pi}}

% Operators and relations
\newunicodechar{×}{\ensuremath{\times}}
\newunicodechar{·}{\ensuremath{\cdot}}
\newunicodechar{±}{\ensuremath{\pm}}
\newunicodechar{−}{\ensuremath{-}}
\newunicodechar{≈}{\ensuremath{\approx}}
\newunicodechar{≤}{\ensuremath{\leq}}
\newunicodechar{≥}{\ensuremath{\geq}}
\newunicodechar{≠}{\ensuremath{\neq}}
\newunicodechar{∈}{\ensuremath{\in}}
\newunicodechar{∂}{\ensuremath{\partial}}
\newunicodechar{∇}{\ensuremath{\nabla}}
\newunicodechar{√}{\ensuremath{\sqrt{}}}
\newunicodechar{∞}{\ensuremath{\infty}}
\newunicodechar{∝}{\ensuremath{\propto}}
\newunicodechar{∑}{\ensuremath{\sum}}
\newunicodechar{∫}{\ensuremath{\int}}
\newunicodechar{ℝ}{\ensuremath{\mathbb{R}}}
\newunicodechar{ℕ}{\ensuremath{\mathbb{N}}}
\newunicodechar{‖}{\ensuremath{\|}}
\newunicodechar{⊕}{\ensuremath{\oplus}}
\newunicodechar{⊗}{\ensuremath{\otimes}}
\newunicodechar{∀}{\ensuremath{\forall}}
\newunicodechar{∃}{\ensuremath{\exists}}
\newunicodechar{¬}{\ensuremath{\neg}}
\newunicodechar{∧}{\ensuremath{\wedge}}
\newunicodechar{∨}{\ensuremath{\vee}}

% Arrows
\newunicodechar{→}{\ensuremath{\rightarrow}}
\newunicodechar{←}{\ensuremath{\leftarrow}}
\newunicodechar{↔}{\ensuremath{\leftrightarrow}}
\newunicodechar{↑}{\ensuremath{\uparrow}}
\newunicodechar{↓}{\ensuremath{\downarrow}}
\newunicodechar{⇒}{\ensuremath{\Rightarrow}}
\newunicodechar{⇐}{\ensuremath{\Leftarrow}}
\newunicodechar{⇔}{\ensuremath{\Leftrightarrow}}

% Super/subscripts written as literal characters
\newunicodechar{²}{\ensuremath{^{2}}}
\newunicodechar{³}{\ensuremath{^{3}}}
\newunicodechar{⁴}{\ensuremath{^{4}}}
\newunicodechar{⁵}{\ensuremath{^{5}}}
\newunicodechar{⁶}{\ensuremath{^{6}}}
\newunicodechar{₀}{\ensuremath{_{0}}}
\newunicodechar{₁}{\ensuremath{_{1}}}
\newunicodechar{₂}{\ensuremath{_{2}}}
\newunicodechar{₃}{\ensuremath{_{3}}}
\newunicodechar{ᵀ}{\ensuremath{^{\mathsf{T}}}}

% Status glyphs. The papers use emoji as status markers in tables
% (implemented / failed / warning). Mapped to typographic equivalents:
% a colour emoji font in lualatex is unreliable, and the MEANING is what
% carries — a reader needs to see "this one failed", not a red cross.
\newunicodechar{✅}{\ensuremath{\checkmark}}
\newunicodechar{✓}{\ensuremath{\checkmark}}
\newunicodechar{❌}{\ensuremath{\times}}
\newunicodechar{✗}{\ensuremath{\times}}
\newunicodechar{⚠}{\textbf{!}}
\newunicodechar{🧩}{\textbf{[piece]}}
\newunicodechar{📋}{\textbf{[list]}}
\newunicodechar{🟡}{\textbf{[partial]}}
\newunicodechar{🔴}{\textbf{[dead]}}
\newunicodechar{🟢}{\textbf{[live]}}
\newunicodechar{▼}{\ensuremath{\blacktriangledown}}
\newunicodechar{►}{\ensuremath{\blacktriangleright}}
\newunicodechar{□}{\ensuremath{\square}}
\newunicodechar{█}{\rule{0.6em}{0.9ex}}
\newunicodechar{•}{\textbullet}
\newunicodechar{†}{\textdagger}
\newunicodechar{‑}{-}
\newunicodechar{‐}{-}
\newunicodechar{…}{\ldots}
\newunicodechar{।}{\textbar}   % danda: rendered by the Bengali fallback in text
\newunicodechar{⁻}{\ensuremath{^{-}}}

% Zero-width formatting codepoints. These are invisible in an editor but
% are real characters to LaTeX: an emoji written as "⚠️" is TWO
% codepoints (U+26A0 U+FE0F), so declaring only the visible one leaves a
% variation selector with no definition and the build dies on a symbol
% nobody can see. Both are mapped to nothing.
\newunicodechar{^^^^fe0f}{}    % VARIATION SELECTOR-16
\newunicodechar{^^^^200d}{}    % ZERO WIDTH JOINER
\newunicodechar{^^^^200b}{}    % ZERO WIDTH SPACE

\usepackage{longtable}
\usepackage{array}

% Comparison matrices and cost ledgers are wide. Let them shrink to the
% text block rather than overflow the margin.
\usepackage{adjustbox}

% Code blocks: wrap long lines instead of running off the page. Several
% papers quote real C++ with long identifiers.
\usepackage{fvextra}
\fvset{breaklines=true,breakanywhere=true,fontsize=\small}

\setlength{\parskip}{0.5em}
\setlength{\parindent}{0pt}

\usepackage{fancyhdr}
\pagestyle{fancy}
\fancyhf{}
\fancyfoot[C]{\thepage}
\fancyhead[L]{\small\itshape Synthetic Intelligence — research lineage}
\fancyhead[R]{\small\itshape irfan@furylogic.com}
\renewcommand{\headrulewidth}{0.2pt}
\setlength{\headheight}{14pt}
\setlength{\headsep}{18pt}
\usepackage{bookmark}
\IfFileExists{xurl.sty}{\usepackage{xurl}}{} % add URL line breaks if available
\urlstyle{same}
\hypersetup{
  pdftitle={Conceptual Mutation Architecture (CMA): Generating Non-Trivial Invariants through Structural Bisociation},
  pdfauthor={Irfan Mahir --- Independent Researcher, Dhaka, Bangladesh --- irfan@furylogic.com},
  colorlinks=true,
  linkcolor={black},
  filecolor={Maroon},
  citecolor={Blue},
  urlcolor={[HTML]{0B5394}},
  pdfcreator={LaTeX via pandoc}}

\title{Conceptual Mutation Architecture (CMA): Generating Non-Trivial
Invariants through Structural Bisociation}
\author{Irfan Mahir --- Independent Researcher, Dhaka, Bangladesh ---
irfan@furylogic.com}
\date{}

\begin{document}
\maketitle

{
\hypersetup{linkcolor=}
\setcounter{tocdepth}{3}
\tableofcontents
}
\section{Conceptual Mutation Architecture (CMA): Generating Non-Trivial
Invariants through Structural
Bisociation}\label{conceptual-mutation-architecture-cma-generating-non-trivial-invariants-through-structural-bisociation}

\textbf{Author:} Irfan Mahir \emph{Independent Researcher, Dhaka,
Bangladesh} irfan@furylogic.com

\begin{center}\rule{0.5\linewidth}{0.5pt}\end{center}

\subsection{Abstract}\label{abstract}

Discrete graph-rewriting engines for Synthetic Intelligence (SI)
traditionally rely on axiomatic deduction and strict topological
constraints. While this ensures high logical coherence and eliminates
the probabilistic ``hallucinations'' characteristic of Large Language
Models (LLMs), it creates a closed-system bottleneck: the engine cannot
generate novelty outside its initial topological training data. We
introduce the Conceptual Mutation Architecture (CMA), an extension for
discrete C++ graph systems that leverages \textbf{Structural
Bisociation}. By systematically injecting forced, non-logical edges
between semantically distant nodes, and subsequently forcing the
engine's inference pipeline to bridge the resulting logical gap, CMA
allows the SI to synthesize original concepts and non-trivial
invariants.

\begin{center}\rule{0.5\linewidth}{0.5pt}\end{center}

\subsection{1. Introduction: The Limits of Deductive
Ontologies}\label{introduction-the-limits-of-deductive-ontologies}

In continuous probabilistic models (LLMs), ``novelty'' is a byproduct of
high-temperature token sampling across a massive, non-linear latent
space. In a discrete graph engine, where concepts are strict nodes and
relationships are explicit structural edges, novelty cannot arise
probabilistically. If an SI engine is constrained to classical deduction
(e.g., Modus Ponens) and induction over a fixed ontology, its thought
processes become purely tautological.

To achieve human-like insight, the engine must perform what Arthur
Koestler termed \emph{Bisociation}---the collision of two mutually
incompatible logical frames. CMA mathematically formalizes this
collision within a graph-theoretical framework.

\subsection{2. Mathematical Framework of the Static
Engine}\label{mathematical-framework-of-the-static-engine}

Let the SI's internal ontology be a directed graph \(G = (V, E)\), where
\(V\) represents the set of concept nodes and \(E\) represents the set
of verified logical relationships (edges).

The semantic distance \(d(v_i, v_j)\) between any two nodes is defined
by the shortest path algorithm across the graph's topology. In the
standard deduction loop, the probability of the engine traversing or
forming a new connection between \(v_i\) and \(v_j\) is inversely
proportional to their distance:

\[ P_{deduce}(v_i \rightarrow v_j) \propto \frac{1}{d(v_i, v_j)^2} \]

This ensures logical safety but absolutely prevents creative leaps. A
node like \texttt{gravity} will never interact with \texttt{RDMA}
because \(d(\text{gravity}, \text{RDMA}) \to \infty\) in the trained
ontology.

\subsection{\texorpdfstring{3. The Mutation Primitive
(\(\mathcal{M}\))}{3. The Mutation Primitive (\textbackslash mathcal\{M\})}}\label{the-mutation-primitive-mathcalm}

CMA introduces a stochastic mutation operator \(\mathcal{M}\) that
periodically overrides the deductive pipeline. Instead of following
existing edges, \(\mathcal{M}\) intentionally targets nodes with maximum
semantic distance to force a hypothesized edge, denoted as \(E_{hyp}\).

The engine selects two nodes \(v_a\) and \(v_b\) such that their
semantic distance exceeds a hyperparameter threshold \(\tau_{novelty}\):

\[ d(v_a, v_b) > \tau_{novelty} \implies E_{hyp}(v_a, v_b) = 1 \]

\subsubsection{3.1 The Bisociation Energy
Function}\label{the-bisociation-energy-function}

Once \(E_{hyp}\) is injected, the engine enters the \textbf{Resolution
Phase}. The system is put under logical ``tension,'' represented by
energy \(U\). The engine must attempt to build a coherent sub-graph path
\(P\) that justifies the connection between \(v_a\) and \(v_b\) using
intermediary nodes.

The energy of the hypothesized connection is defined as:

\[ U_{hyp} = \sum_{k=1}^{n} R(e_k) - \alpha \cdot C(v_a, v_b) \]

Where:

\begin{itemize}
\tightlist
\item
  \(R(e_k)\) is the logical resistance of the intermediary edges used to
  build the bridge.
\item
  \(C(v_a, v_b)\) is the structural conductivity of the new conceptual
  bridge.
\item
  \(\alpha\) is a scaling constant representing the system's current
  ``exploratory state'' (analogous to dopamine).
\end{itemize}

\subsection{4. Logical Resolution and The Novelty
Test}\label{logical-resolution-and-the-novelty-test}

The forced edge \(E_{hyp}\) is initially treated as a logical
contradiction. The SI's core engine must resolve this contradiction by
generating intermediate analogies.

For example, if \(\mathcal{M}\) forces
\(E_{hyp}(\text{gravity}, \text{RDMA})\):

\begin{enumerate}
\def\labelenumi{\arabic{enumi}.}
\tightlist
\item
  The engine decomposes \texttt{gravity}:
  \(\text{gravity} \to \text{attractive\_force} \to \text{data\_pull}\).
\item
  The engine decomposes \texttt{RDMA}:
  \(\text{RDMA} \to \text{direct\_memory\_access} \to \text{data\_pull}\).
\item
  The bridge is formed at the intersection node \texttt{data\_pull}.
\end{enumerate}

If the engine successfully constructs a path \(P^*\) that minimizes the
bisociation energy \(U_{hyp}\) below a critical threshold \(U_{crit}\),
the edge is validated.

\[ \text{If } \min(U_{hyp}) < U_{crit} \implies E_{hyp} \to E_{verified} \]

The graph has now permanently adopted a novel logical connection that
was not in its initial training set. The SI has ``invented'' a new idea.

\subsection{5. Algorithmic Pruning
(Decay)}\label{algorithmic-pruning-decay}

Not all mutations result in profound insights; many are simply
nonsensical. If the resolution pipeline fails to find a path \(P^*\)
that satisfies the energy threshold within \(N\) clock cycles, the
hypothesized edge undergoes rapid topological decay:

\[ \frac{d}{dt} E_{hyp} = -\lambda E_{hyp} \]

Where \(\lambda\) is a severe decay constant. The edge is pruned from
the graph, preventing the graph from filling with unresolved,
contradictory noise.

\subsection{6. Conclusion}\label{conclusion}

The Conceptual Mutation Architecture liberates discrete graph-rewriting
systems from the trap of tautological deduction. By formalizing
Koestler's bisociation as a stochastic topological operator
\(\mathcal{M}\) constrained by an energy-minimization resolution phase,
the SI achieves true structural novelty. It gains the ability to
``brainstorm'' by forcing disparate concepts together and rigorously
testing their logical intersections, bridging the gap between static
knowledge retrieval and dynamic human-like creativity.

\end{document}
